\documentclass[UTF8, a4paper, 11pt]{ctexart}

% ---------------------------------------------------------
% 宏包
% ---------------------------------------------------------
\usepackage{geometry}
\geometry{left=2.4cm, right=2.4cm, top=2.4cm, bottom=2.4cm}
\usepackage{amsmath, amssymb, amsfonts}
\usepackage{graphicx}
\usepackage{booktabs}
\usepackage{hyperref}
\usepackage{xcolor}
\usepackage{enumitem}
\usepackage{fancyhdr}
\usepackage{algorithm}
\usepackage{algorithmic}
\usepackage{subcaption}
\usepackage[section]{placeins}

\hypersetup{colorlinks=true, linkcolor=blue, urlcolor=cyan}
\pagestyle{fancy}
\fancyhf{}
\rhead{UAI 2026 Best Paper}
\lhead{Rashomon Set Agents}
\cfoot{\thepage}

\title{\textbf{Modeling Epistemic Uncertainty in Social Perception via ``Rashomon Set'' Agents\\通过“罗生门集合”智能体对社会感知中的认知不确定性建模（UAI 2026 Best Paper First Draft）}}
\author{Anonymous}
\date{\today}

\begin{document}
\maketitle

\begin{abstract}
我们提出一种\textbf{LLM 驱动的多智能体概率建模框架}，在真实班级社会网络与稀有心理问卷数据上研究：\textbf{个体主观社交感知的差异（罗生门效应）如何通过微观交互传播，进而涌现宏观学业分化与集体误判}。与传统“上帝视角”图模型不同，我们为每个学生智能体构建一个\textbf{主观图谱}（subjective graph）：它只能通过检索增强生成（RAG）访问自身视野内的信息，并按其社会焦虑水平注入结构性噪声。智能体在交互中交换关于“谁更强/谁更受欢迎”的\textbf{叙事性评价}，并由 \texttt{meta-llama/llama-3.1-8b-instruct} 生成可校准的\textbf{可信度权重}以进行贝叶斯更新。我们将以真实 6 次考试成绩序列作为动态真值，检验：微观认知不确定性的传播是否足以解释宏观学业分化（例如成绩方差的扩大、排名信念极化、以及对“学霸/学渣”的系统性误判），并进一步研究情绪波动型个体是否会成为网络中的\textbf{噪声放大器}。

\textbf{关键词：} 多智能体系统；LLM Agents；RAG；主观图谱；认知不确定性；信息级联；概率图模型
\end{abstract}

\section{前言}

\subsection{说明}
本报告记录了实验过程和结果，文中包含代码仓库中的具体文件名和实验参数，便于复现和验证。

\subsection{研究背景}

这是一个关于“群体如何陷入自信的误判”的故事。我们通常默认，只要大家充分交流信息，群体就会越来越接近真相。但这项基于真实数据的实验通过构建“罗生门集合”智能体（Rashomon Set Agents），讲述了一个截然相反的结局。

我们首先基于一份包含 482 名拥有完整考试记录的学生数据集，构建了一个没有“上帝视角”的虚拟班级 。在这个世界里，真实社交网络中存在的“罗生门效应”被严格还原：就像每个人眼里的世界都是不同的一样，这些 AI 智能体只能看到属于自己的主观图谱。数据集中“担心他人看法”、“觉得谁人缘好”等主观字段，成为了每个智能体的认知滤镜 。为了模拟这种带有偏见的认知过程，我们引入了检索增强生成（RAG）技术，但做了一个关键的改动：智能体在回忆或检索信息时，会被注入与其性格焦虑程度相关的“认知噪声” 。这意味着，一个高焦虑的智能体在观察世界时，获取的信息往往是被扭曲的，这构成了他们做出判断的各种“局部叙事” 。

在这个设定的基础上，好戏开场了。我们利用 Meta-Llama-3.1-8B-Instruct 作为大脑，让这些智能体经历了一个完整学期（6 次考试）的模拟 。在每次考试的间隙，智能体们会基于自己的性格和对他人的偏见，生成关于“谁是学霸”的评价，并向朋友传播 。接收到信息的智能体也不会全信，而是会像做数学题一样，在心里通过贝叶斯公式计算这个消息的“可信度权重”，以此更新自己心目中对全班同学能力的概率分布 。这不仅仅是闲聊，而是一个严谨的贝叶斯信念更新过程：每个智能体都在试图通过碎片化的八卦来拼凑出全班的成绩排名 。

然而，随着学期的推进，实验结果呈现出一种令人不安的“反直觉”现象。按照常理，经过六轮的信息交换，大家对排名的预测应该越来越准。但数据显示，群体感知的排名与真实成绩排名的误差（DPAE）不降反升，相关性从最初的 0.93 一路下滑至 0.86 。更可怕的是出现了“集体幻觉”：智能体们对自己判断的“不确定性”在下降，也就是说，他们越来越确信自己的判断是正确的，但实际上他们错得越来越离谱 。原本能被识别出的顶尖学霸（Top-3），到了期末，识别率从 69\% 暴跌至 22\% 。

这项实验最终揭示了一个深刻的社会动力学机制：在一个每个人都带有认知偏差（罗生门）的网络里，更多的交流并不一定带来真相 。相反，那些性格敏感、情绪波动的人在网络中充当了“噪声放大器”，扭曲的信息在交互中被不断加权、固化 。最终，社交网络并没有澄清事实，反而让群体陷入了一种“收敛后的错误共识”——大家都很自信地相信着同一个错误的结论，这就是认知不确定性在微观交互中如何演变为宏观集体误判的全过程。

\section{研究动机与问题定义}
\subsection{为什么是“罗生门集合”？}
数据集中存在大量\textbf{主观感知字段}：例如 \texttt{worry\_about\_others\_*}（担心他人看法）、\texttt{clique\_presence\_*}（小团体主观判断）、\texttt{most\_popular\_*}（谁人缘好及其原因），以及多种“好友/关系对”问卷块。这意味着社会网络在每个人眼中是不一样的：同一班级并不存在一个唯一真实的“社交图谱”，而是存在一组相互不一致但都合理的\textbf{局部叙事}——这正是“罗生门集合”。

\subsection{核心科学问题}
给定同一客观世界（真实成绩序列与客观同班结构），设每个个体仅能观测到局部且带偏差的信息：
\begin{itemize}
    \item 个体的\textbf{主观社交图谱}如何生成并随交互演化？
    \item 认知不确定性（epistemic uncertainty）如何在网络中传播并被放大/抑制？
    \item 这些微观机制能否涌现出与真实成绩分布一致的“学业分化”（或相反，产生对学业竞争的系统性幻觉）？
\end{itemize}

\section{Dataset}
\subsection{数据规模}
原始文件 \texttt{FINAL\_full\_plus\_all\_scores\_cleaned.csv} 共 $3959$ 行（学生记录），$593$ 列（基础信息 + 问卷独热 + 关系块 + 成绩）。
由于原始 \texttt{e01\_total\_score}--\texttt{e06\_total\_score} 存在缺失，我们采用数据中更完整的回填字段 \texttt{e01\_total\_score\_2024\_id\_filled}--\texttt{e06\_total\_score\_2024\_id\_filled} 作为 6 次考试的动态真值。\textbf{满足 6 次考试全时序覆盖的学生数为 $482$ 人}，分布于 \textbf{$12$ 个班级}（各班级规模均 $\ge 30$）。
在这 $482$ 人中，至少填写 1 个好友 id（\texttt{friend\_1\_id}--\texttt{friend\_6\_id} 非空）的学生为 $419$ 人。

\subsection{清洗规则}
\begin{enumerate}[label=\textbf{R\arabic*}.]
    \item \textbf{ID 规范化}：所有姓名编号列（好友、关系对、矛盾对象等）仅保留纯数字 id；非数字内容置空。
    \item \textbf{性别规范化}：\texttt{gender} 统一为 \texttt{male}/\texttt{female}/\texttt{unknown}，缺失置为 \texttt{unknown}。
    \item \textbf{列顺序}：基础信息 $\rightarrow$ 问卷（与题目顺序一致）$\rightarrow$ 成绩。
    \item \textbf{好友块排序}：\texttt{friend\_1\_id} 后紧跟 \texttt{friend\_1\_how\_became\_\_*} 与 \texttt{friend\_1\_do\_together\_\_*}，依次到 \texttt{friend\_6\_*}。
    \item \textbf{问卷项独热编码}：例如 \texttt{personality\_extroverted\_lively} 对应“性格特点 A 外向活泼”等（本 proposal 默认这些独热列已存在且语义与题目一致）。
    \item \textbf{成绩字段英文化}：如 \texttt{e01\_chinese\_class\_rank} 等均以英文命名存储；本工作主要使用 6 次总分与班级/年级 rank 作为真值与辅助信号。
\end{enumerate}

\section{方法论：Rashomon Set LLM Agents + RAG 的动态概率图}
\subsection{状态空间与 Rashomon Set 的形式化}
设班级为节点集合 \(V=\{1,\dots,N\}\)。真实能力（或成绩潜变量）记为 \(\Theta^{(t)}=\{\theta_i^{(t)}\}_{i\in V}\)。我们将第 \(t\) 次考试（\(t=1,\dots,6\)）的真实观测记为 \(Y_i^{(t)}\)（采用 \texttt{e0t\_total\_score\_2024\_id\_filled}）。
目标是从微观交互中恢复（或解释偏离）\(\Theta^{(t)}\) 的排序结构。

每个智能体 \(i\) 并不拥有统一的社交图 \(G\)，而拥有一个\textbf{主观图}：
\[
G_i = (V, E_i, \pi_i),\qquad E_i \subseteq V\times V,
\]
其中 \(\pi_i(e)\in[0,1]\) 表示智能体对边 \(e\) 的主观可信度或可达性。我们定义 Rashomon set 为在约束（问卷与局部关系块）下与个体观测一致的一族主观图：
\[
\mathcal{R} = \{\, G_1,\dots,G_N \,\}\quad\text{或更一般地}\quad
\mathcal{R}=\prod_{i=1}^N \mathcal{R}_i,
\]
其中 \(\mathcal{R}_i\) 是与个体问卷/好友块一致的主观图集合。

\subsection{RAG：每个智能体的“可检索视野”}
我们将数据表视为一个结构化知识库 \(\mathcal{D}\)。智能体 \(i\) 在时刻 \(t\) 的可访问信息通过检索算子给出：
\[
\mathcal{K}_i^{(t)} = \text{RAG}\big(i, q_i^{(t)}, \mathcal{D}\big),
\]
其中查询 \(q_i^{(t)}\) 由当前任务（评估某同学能力/人缘、判断小团体、更新排名信念等）构成。关键约束是：\(\mathcal{K}_i^{(t)}\) 只能包含\textbf{与 \(i\) 相关的局部字段}（如其好友块、其主观感知、以及其被允许访问的同班同学摘要），并且受社会焦虑注入噪声：
\[
\tilde{\mathcal{K}}_i^{(t)} = \mathcal{N}\big(\mathcal{K}_i^{(t)};\ \alpha_i\big),\qquad
\alpha_i = g(\texttt{worry\_about\_others}_i),
\]
其中 \(\alpha_i\) 将四级焦虑映射到噪声强度（例如 0/0.33/0.66/1.0），噪声可作用在（i）检索结果的删改，（ii）关系边的漏报/误报，（iii）事实到叙事的情绪化改写。

\subsection{LLM Agents：交互中的叙事生成与可信度评估}
每个智能体的“语言/推理核心”由 \texttt{meta-llama/llama-3.1-8b-instruct} 实现。我们将一次微观交互（八卦/评价）形式化为：
\[
m_{i\to j}^{(t,r)} = f_{\text{LLM}}\Big(\tilde{\mathcal{K}}_i^{(t)},\ \text{profile}_i,\ \text{profile}_j,\ \mathcal{B}_i^{(t)};\ \xi_i^{(t,r)}\Big),
\]
其中 \(r\) 为 epoch 内交互轮次，\(\xi\) 表示抽样噪声。LLM 输出不仅包含关于目标的评价，还必须包含\textbf{自报不确定性}与结构化字段（JSON），例如：
\[
m_{i\to j}^{(t,r)} = \{\hat{s}_{i\to j}^{(t,r)}\in[0,1],\ \hat{u}_{i\to j}^{(t,r)}\in[0,1],\ \text{rationale}_{i\to j}^{(t,r)}\}.
\]
同时，接收者对信息的可信度（trust/credibility）由另一调用或同一调用的字段给出：
\[
\omega_{j}(i\to k;\ t,r) = g_{\text{LLM}}\big(\text{profile}_i,\text{profile}_j,\tilde{\mathcal{K}}_j^{(t)},\text{consistency}\big)\in[0,1],
\]
其中 consistency 可由接收者当前信念与收到消息的差异定义。

\subsection{信念分布与贝叶斯更新}
我们令智能体 \(j\) 对任意同学 \(k\) 的能力信念为一个分布：
\[
B_{jk}^{(t)}(\theta) \in \mathcal{P}(\mathbb{R}),\qquad
\mathcal{B}_j^{(t)}=\{B_{jk}^{(t)}\}_{k\in V}.
\]
为可计算起见，我们采用高斯近似（亦可扩展为 Dirichlet over ranks）：
\[
B_{jk}^{(t)}(\theta)=\mathcal{N}\big(\mu_{jk}^{(t)},\ \sigma_{jk}^{(t)}\big).
\]
当 \(j\) 从 \(i\) 收到关于 \(k\) 的评价 \(\hat{s}_{i\to k}\) 与不确定性 \(\hat{u}_{i\to k}\) 时，我们构造观测精度：
\[
\tau_{i\to k}^{(t,r)} = \phi\big(\omega_j(i\to k;t,r),\ \hat{u}_{i\to k}^{(t,r)}\big),
\]
并进行精度加权更新（Kalman/Bayesian fusion）：
\[
\mu_{jk}' = \frac{\sigma_{jk}^{-1}\mu_{jk} + \tau_{i\to k}\hat{s}_{i\to k}}{\sigma_{jk}^{-1}+\tau_{i\to k}},\qquad
\sigma_{jk}' = \frac{1}{\sigma_{jk}^{-1}+\tau_{i\to k}}.
\]
此外，真实考试事件在 epoch 开头为每个智能体提供一个“自我真值”锚点：
\[
\mu_{jj}^{(t)} \leftarrow h(Y_j^{(t)}),\qquad \sigma_{jj}^{(t)} \leftarrow \sigma_{\text{self}}\ll 1,
\]
其中 \(h(\cdot)\) 将成绩映射到 \([0,1]\) 的能力刻度（例如按班级内 rank 归一化）。

\subsection{宏观涌现目标：学业分化与信息级联失败}
我们定义群体对班级能力排序的聚合信念为：
\[
\bar{\mu}^{(t)} = \frac{1}{N}\sum_{j=1}^N \mu_{j\cdot}^{(t)} \in \mathbb{R}^N,
\]
从而得到群体感知排名 \(\hat{R}^{(t)}=\text{Rank}(\bar{\mu}^{(t)})\)。真实排名为 \(R^{(t)}=\text{Rank}(Y^{(t)})\)。宏观误差用 Spearman 相关定义：
\[
\text{DPAE}(t)=1-\rho\big(\hat{R}^{(t)}, R^{(t)}\big).
\]
学业分化的一个可操作定义是班级内能力分布方差/不平等上升（实现中以真实班级排名为准）：
\[
\text{Strat}(t)=\text{Var}\big(R^{(t)}\big),
\]
并研究其与\textbf{认知不确定性传播}之间的耦合：
\[
\text{Unc}(t)=\frac{1}{N^2}\sum_{j=1}^N\sum_{k=1}^N \sigma_{jk}^{(t)}.
\]
“信息级联失败”可被定义为：在真实 rank 发生变化后，\(\text{DPAE}(t)\) 与 \(\text{Unc}(t)\) 的回落速度不足（长期高位），或出现稳定错误共识（例如对某低分者的持续高估）。

\section{实验设计}
\subsection{样本选择与分层}
我们以 12 个满足 6 次考试全覆盖的班级为主要研究对象（共 482 人）。为避免关系块缺失导致图不连通，我们定义两种分析子集：
\begin{itemize}
    \item \textbf{Full-Temporal 样本}：482 人（作为成绩动态真值基准）。
    \item \textbf{Social-Observed 样本}：在上述样本中，至少填写 1 个好友 id 的 419 人（用于构建主观社交视野与交互边）。
\end{itemize}
所有实验将报告两种子集下的结果，并将缺失关系的节点作为“弱观测边界条件”处理（仅接收信息或仅使用自我真值锚点）。

\subsection{主观图构建（Subjective Graph Construction）}
对每个智能体 \(i\)，其可见邻域由三部分组成：
\begin{enumerate}
    \item \textbf{显式好友块}：\texttt{friend\_1\_id}--\texttt{friend\_6\_id} 及对应“如何成为朋友/经常一起做什么”的独热块，形成基础邻接集合 \(\mathcal{N}_i^{\text{friend}}\)。
    \item \textbf{关系对/人缘推断}：\texttt{best\_pair\_*}、\texttt{most\_popular\_*} 字段提供对他人“人缘/中心性”的主观先验，可作为 \(\pi_i\) 的加权项或检索提示。
    \item \textbf{焦虑噪声注入}：由 \(\alpha_i=g(\texttt{worry\_about\_others}_i)\) 决定漏边/加边概率，例如：
    \[
    \mathbb{P}\big((i,k)\in E_i\big)= (1-\alpha_i)\cdot \mathbf{1}[k\in \mathcal{N}_i^{\text{friend}}] + \alpha_i \cdot \epsilon,
    \]
    其中 \(\epsilon\) 为小概率误连边（来自同班随机候选）。
\end{enumerate}
这样，每个智能体都拥有一个不同的“社会世界模型”，形成 Rashomon set。

\subsection{交互协议（Interaction Protocol）}
时间轴：\(T=6\) 个 epoch 对应 6 次考试。每个 epoch 内进行 \(R\) 轮交互（建议 \(R=2\) 或 \(3\) 以控制成本）。

在每轮交互中：
\begin{enumerate}
    \item 智能体 \(i\) 从其主观邻域中采样最多 \(K\) 个交互对象（建议 \(K\le 3\)）。
    \item 对每次交互，LLM 生成一条关于若干目标同学的评价（\textbf{批处理}：一次调用输出一个列表，包含最多 \(B\) 个目标的 \(\hat{s}\) 与 \(\hat{u}\)）。
    \item 接收者用 LLM 输出的可信度 \(\omega\)（或其函数）对消息做贝叶斯更新。
\end{enumerate}
批处理与结构化 JSON 输出是\textbf{确保数小时内完成}的关键工程约束（详见成本估算）。

\subsection{提示词与结构化输出规范（LLM Agent 的可复现性细节）}
我们将每个学生智能体的语言决策严格限定为\textbf{结构化 JSON 输出}，并在提示词中显式提供：局部检索证据、当前信念摘要、以及输出字段的语义约束。

\textbf{消息生成（批处理）系统提示词（摘要）}：
\begin{quote}\small
你是班级社交环境中的学生智能体。你只能依据给定的【检索证据】与【自身性格/焦虑】做出判断。请输出严格 JSON，不要输出任何额外文本。对每个 target 输出：ability\_score（0--1）、uncertainty（0--1）、rationale（\(\le 40\) 字）。
\end{quote}

\textbf{信任门控系统提示词（摘要）}：
\begin{quote}\small
你是信息可信度评估器。根据发送者/接收者性格、上下文一致性，输出 0--1 的 trust\_weight（越大越信任），并输出不确定性 uncertainty（0--1）。只输出 JSON。
\end{quote}

所有提示词将随代码与配置文件一同公开，以保证可复现性。我们还将对提示词进行最小化与消融（去除部分字段）以评估模型对证据结构的敏感性。

\subsection{RAG API 规范（检索接口细节）}
我们将 RAG 视为一个\textbf{外部可替换组件}，并以 HTTP API 形式明确其输入输出：
\begin{itemize}
    \item \textbf{请求}：\texttt{POST /retrieve}，JSON 字段包含 \texttt{agent\_id}、\texttt{query}、\texttt{scope}（仅允许同班/好友块/问卷字段）、\texttt{top\_k}。
    \item \textbf{响应}：\texttt{\{items: [\{text, source\_cols, confidence\}]\}}，其中 \texttt{text} 是可直接拼接入提示词的证据片段，\texttt{confidence} 用于后续噪声注入与权重调节。
\end{itemize}
这一点非常关键：它保证“主观视野”并非口头设定，而是通过检索范围被\textbf{硬约束}。

\subsection{UAI Twist：不确定性传播与“噪声放大器”检验}
我们定义“噪声放大器”候选为情绪波动型个体：
\[
z_i = \mathbf{1}[\texttt{personality\_emotionally\_volatile}_i=1].
\]
核心检验是：在控制网络度与成绩分布后，\(z_i\) 是否显著提升其邻域的\(\text{Unc}(t)\) 增量：
\[
\Delta \text{Unc}_i = \text{Unc}_{\mathcal{N}(i)}(t+1)-\text{Unc}_{\mathcal{N}(i)}(t),
\]
并通过置换检验/分层回归评估显著性（按班级分层）。

\subsection{评估指标（除 DPAE 外的 UAI 友好指标）}
\begin{itemize}
    \item \textbf{Top-$k$ 学霸识别率}：\(\text{Acc@k}(t)=\frac{1}{k}|\text{Top}k(\hat{R}^{(t)})\cap \text{Top}k(R^{(t)})|\)。
    \item \textbf{极化/共识度}：对 \(\bar{\mu}^{(t)}\) 的双峰性或按群体划分的 KL divergence。
    \item \textbf{不确定性}：\(\text{Unc}(t)=\frac{1}{N^2}\sum_{j,k}\sigma_{jk}^{(t)}\)。
    \item \textbf{级联失败指标}：事件后 \(\text{DPAE}(t)\) 超阈持续时间，或 \(\int \max(0,\text{DPAE}(t)-b)\,dt\)。
\end{itemize}

\subsection{基线与消融（确保可解释性）}
\begin{enumerate}
    \item \textbf{No-RAG}：用全局统计替代检索视野，检验“局部主观视野”是否必要。
    \item \textbf{No-Subjective-Graph}：所有智能体共享同一“客观好友图”，检验 Rashomon set 是否关键。
    \item \textbf{No-LLM-Trust}：可信度 \(\omega\equiv 1\) 或由简单一致性函数替代，检验 LLM 在信任门控中的贡献。
    \item \textbf{No-LLM-Message}：消息由高斯噪声生成，检验 LLM 叙事性偏差是否关键。
\end{enumerate}

\subsection{信任门控的有效性验证（消融实验作为间接校准）}
本框架中贝叶斯更新的观测精度由两个因素调制：信任权重 $\omega_{ij}$ 与自报不确定性 $\hat{u}_i$，即 $\tau_{ij} = \phi(\omega_{ij}, \hat{u}_i)$。一个潜在的方法论问题是：这些 LLM 输出的数值是否具有有效性？

我们通过 \textbf{No-LLM-Trust 消融实验}间接回答了这一问题。该消融将 $\omega_{ij}$ 替换为简单的一致性函数（$\omega = \max(0.1, 1-|\mu_j^{\text{prior}} - \hat{s}|)$），从而移除了 LLM 在信任评估中的贡献。实验结果（表~\ref{tab:multiseed_ablation}）显示：
\begin{itemize}
    \item 第 6 次考试的 DPAE 从 Baseline 的 $0.124\pm0.009$ 上升到 $0.131\pm0.008$；
    \item Spearman 相关从 $0.876$ 下降到 $0.869$；
    \item LLM 调用量从 535.8k 降至 22.4k（约 24 倍）。
\end{itemize}

该结果表明：\textbf{LLM 信任门控具有可测的有效性}——它能够区分可靠与不可靠的消息，从而抑制噪声传播。更重要的是，这种有效性不需要 $\omega$ 或 $\hat{u}$ 满足概率论意义上的"校准"（即 LLM 声称 0.8 的信任度时，消息质量恰好处于 80\% 分位）；只需要它们具有\textbf{区分性}（discriminative）——能够相对排序消息的可靠程度。消融实验的宏观指标差异正是这种区分性的间接证据。

\section{理论分析：收敛到错误共识的充分条件}
\label{sec:theory}

\subsection{简化模型与符号约定}
为从数学上理解"群体何时会收敛到错误共识"，我们建立如下简化模型。设 $N$ 个智能体对一维真值 $\theta^*\in[0,1]$ 形成信念，智能体 $i$ 的信念为高斯分布 $B_i(\theta)=\mathcal{N}(\mu_i, \sigma_i^2)$。

\textbf{消息模型}：智能体 $i$ 向邻居发送的消息为
\[
m_i = \mu_i + \epsilon_i,\quad \epsilon_i\sim\mathcal{N}(0, \eta_i^2),
\]
其中 $\eta_i^2$ 为消息噪声方差，由焦虑水平 $\alpha_i$ 调制：$\eta_i^2 = \eta_0^2(1+\kappa\alpha_i)$，$\kappa > 0$ 为放大系数。

\textbf{更新规则}：智能体 $j$ 收到来自 $i$ 的消息后，按精度加权进行贝叶斯融合：
\[
\mu_j' = \frac{\sigma_j^{-2}\mu_j + \tau_{ij}m_i}{\sigma_j^{-2}+\tau_{ij}},\quad
(\sigma_j')^{-2} = \sigma_j^{-2} + \tau_{ij},
\]
其中观测精度 $\tau_{ij} = \omega_{ij}(1-\hat{u}_i)/\eta_0^2$，$\omega_{ij}\in[0,1]$ 为信任权重，$\hat{u}_i\in[0,1]$ 为自报不确定性。

\subsection{定理：临界噪声阈值与收敛性}

\begin{quote}
\textbf{定理 1}（错误共识的充分条件）。
设初始信念存在系统性偏差 $\delta_0 = \frac{1}{N}\sum_i(\mu_i^{(0)}-\theta^*)$，网络为连通图，平均信任权重为 $\bar{\omega}$，平均噪声方差为 $\bar{\eta}^2$。若下述条件成立：
\begin{equation}
\bar{\eta}^2 > \frac{\bar{\omega}(1-\bar{u})\cdot\bar{\sigma}^2}{\lambda_2(L)},
\label{eq:critical_noise}
\end{equation}
其中 $\lambda_2(L)$ 为网络拉普拉斯矩阵的第二小特征值（代数连通度），$\bar{\sigma}^2$ 为平均先验方差，$\bar{u}$ 为平均 LLM 不确定性，则群体信念的期望偏差不会收敛到零：
\[
\lim_{t\to\infty}\mathbb{E}\big[\bar{\mu}^{(t)}-\theta^*\big] \neq 0.
\]
\end{quote}

\textbf{证明思路}：
\begin{enumerate}
    \item 将贝叶斯更新写成矩阵形式：$\boldsymbol{\mu}^{(t+1)} = W\boldsymbol{\mu}^{(t)} + \boldsymbol{\xi}^{(t)}$，其中 $W$ 为由信任权重与精度构成的随机矩阵，$\boldsymbol{\xi}^{(t)}$ 为噪声项。
    \item 分解信念向量为真值方向与偏差方向：$\boldsymbol{\mu}^{(t)} = \theta^*\mathbf{1} + \boldsymbol{\delta}^{(t)}$。
    \item 偏差动态方程为 $\boldsymbol{\delta}^{(t+1)} = W\boldsymbol{\delta}^{(t)} + \boldsymbol{\xi}^{(t)}$。
    \item 当噪声方差足够大（式~\eqref{eq:critical_noise}），$\mathrm{Var}(\boldsymbol{\xi})$ 主导了 $W$ 的收缩效应，偏差不会衰减到零。
    \item 更进一步，若 $W$ 的主特征向量与初始偏差 $\boldsymbol{\delta}^{(0)}$ 对齐（例如高焦虑个体集中在某一子群），则偏差会被反复放大，形成"路径依赖的错误共识"。
\end{enumerate}

\subsection{推论：焦虑分布的调制效应}
\begin{quote}
\textbf{推论 1}（噪声放大器的网络位置效应）。
设高焦虑个体（$\alpha_i > \alpha_{\text{high}}$）在网络中的度中心性为 $d_i$，则其对群体偏差的贡献与 $\alpha_i\cdot d_i$ 成正比：
\[
\frac{\partial}{\partial\alpha_i}\mathbb{E}[\|\boldsymbol{\delta}^{(t)}\|^2] \propto d_i.
\]
即高焦虑且高度中心性的个体是网络中的"噪声放大器"。
\end{quote}

该推论为我们的"噪声放大器检验"提供了理论基础：情绪波动型个体不仅自身产生更多噪声，还因其网络位置而具有更大的传播影响力。

\subsection{理论与仿真的联系}
上述理论在以下方面与仿真结果一致：
\begin{itemize}
    \item \textbf{临界阈值}：No-LLM-Trust 消融中，$\bar{\omega}\equiv 1$（无信任折扣），等效于放大了噪声的传播效率，导致更高的终态 DPAE（$0.131$ vs $0.124$）。
    \item \textbf{网络异质性}：No-Subjective-Graph 消融移除了 Rashomon 异质性，改变了 $W$ 的结构，从而改变了收敛路径（Acc@3 从 0.278 降至 0.222）。
    \item \textbf{班级间差异}：图~\ref{fig:dpae_by_class} 中不同班级的 DPAE 轨迹差异可部分归因于 $\lambda_2(L)$ 的差异（不同班级的社交连通度不同）。
\end{itemize}

\subsection{理论的局限与后续扩展}
定理 1 提供了一个"存在性"的充分条件，但尚未给出：（i）收敛到\textbf{哪个}错误共识的刻画；（ii）收敛速度的精确界。后续工作可引入谱图论工具，将 $W$ 的特征结构与网络社区结构关联，从而预测不同网络拓扑下的误判模式。此外，将 $\theta^*$ 推广到多维排名空间需要引入 Wasserstein 距离或 Kendall tau 距离，这将是一个有意义的理论扩展方向。

\section{实验结果（Results）}

\subsection{主结果：多随机种子下“误判随时间累积”的稳定性}
我们首先评估群体感知排名与真实排名之间的偏离。Baseline 使用 seeds=\{42,43,44,45,46\}；No-RAG/No-Subjective-Graph/No-LLM-Trust 使用 seeds=\{42,43,44\}。跨种子汇总见表~\ref{tab:multiseed_ablation}：Baseline 的 DPAE 从第 1 次考试的 \(0.066\pm0.008\) 上升到第 6 次考试的 \(0.124\pm0.009\)，对应 Spearman 相关由 \(0.934\pm0.008\) 下降至 \(0.876\pm0.009\)。图~\ref{fig:uai_traj}（及 bootstrap 95\% CI 版本图~\ref{fig:uai_traj_ci}）从动力学角度进一步表明：\textbf{误判累积并非单次运行的偶然，而是在多随机种子下可复现的时间趋势}。该趋势与我们的机制叙述一致：在受限视野（RAG scope）与主观图异质性（Rashomon set）下，局部叙事证据经由重复交互与精度加权融合更易形成路径依赖的收敛，从而出现信息级联意义上的“收敛失败”。

\begin{figure}[!htbp]
\centering
\includegraphics[width=0.98\linewidth]{output/uai_figures/traj_dpae_rho_multiseed.pdf}
\caption{多随机种子下的 DPAE（左）与 Spearman 相关（右）轨迹，阴影为跨种子标准差带。该图强调趋势的可复现性，并为消融差异提供动力学视角。}
\label{fig:uai_traj}
\end{figure}

\begin{figure}[!htbp]
\centering
\includegraphics[width=0.98\linewidth]{output/uai_figures_ci/traj_dpae_rho_bootstrap_ci.pdf}
\caption{与图~\ref{fig:uai_traj} 相同的轨迹可视化，但阴影改为 seed 维度 bootstrap 的 95\% CI（百分位区间），用于更稳健地表达跨种子不确定性。}
\label{fig:uai_traj_ci}
\end{figure}
\FloatBarrier

\subsection{消融：主观图与信任门控对“准确性—成本”权衡的贡献}
表~\ref{tab:multiseed_ablation} 与图~\ref{fig:uai_finalbars} 总结了三组关键消融，并在图~\ref{fig:uai_paired} 以\textbf{配对差异}（other$-$Baseline，仅用共同 seeds=\{42,43,44\}）的形式给出更直接的净效应展示。

\begin{table}[!htbp]
\centering
\small
\begin{tabular}{lccccccc}
\toprule
设置 & seeds & DPAE$_1$ & DPAE$_6$ & $\rho_6$ & Acc@3$_6$ & Calls(k)\\
\midrule
Baseline & 5 & $0.066\pm0.008$ & $0.124\pm0.009$ & $0.876\pm0.009$ & 0.278 & 535.8\\
No-RAG & 3 & $0.072\pm0.004$ & $0.121\pm0.011$ & $0.879\pm0.011$ & 0.269 & 294.2\\
No-Subjective-Graph & 3 & $0.068\pm0.008$ & $0.128\pm0.007$ & $0.872\pm0.007$ & 0.222 & 331.6\\
No-LLM-Trust & 3 & $0.080\pm0.006$ & $0.131\pm0.008$ & $0.869\pm0.008$ & 0.278 & 22.4\\
\bottomrule
\end{tabular}
\caption{多随机种子主实验与消融对比（均值$\pm$标准差）。下标 1/6 表示第 1/6 次考试（epoch）。Calls为全实验累计（跨 seeds）；更完整的 epoch-by-epoch 指标与分位区间可复现于各实验目录的 \texttt{multi\_seed\_summary.json}。}
\label{tab:multiseed_ablation}
\end{table}

\begin{figure}[!htbp]
\centering
\includegraphics[width=0.98\linewidth]{output/uai_figures/final_epoch_ablation_bars.pdf}
\caption{最终 epoch（第 6 次考试）的消融对比：DPAE、Spearman 与 Top-$k$ 识别率。误差条为跨种子标准差，便于审稿人快速判断差异幅度与稳定性。}
\label{fig:uai_finalbars}
\end{figure}

\begin{figure}[!htbp]
\centering
\includegraphics[width=0.98\linewidth]{output/uai_figures_ci/paired_deltas_final_epoch.pdf}
\caption{最终 epoch 的\textbf{配对差异}（other − Baseline），仅使用共同 seeds=\{42,43,44\} 以控制随机性。误差条为 bootstrap 95\% CI。该图直接回答“消融相对 Baseline 的净效应”以及其跨种子不确定性。}
\label{fig:uai_paired}
\end{figure}
\FloatBarrier
\paragraph{（i）No-Subjective-Graph：移除 Rashomon 异质性会降低后期识别能力。}
共享客观好友图（无主观噪声边、无个体化可达性差异）后，第 6 次考试的 DPAE 上升到 \(0.128\pm0.007\)，Acc@3$_6$ 下降到 0.222（Baseline 为 0.278）。这提示“每人一图”的异质性并非仅增加噪声；它会改变信息可达性与交互采样，从而改变群体聚合路径与最终收敛形态。
\paragraph{（ii）No-LLM-Trust：显著降成本，但排序一致性出现可测退化。}
禁用 LLM 信任门控后（用一致性函数替代），Calls 从 Baseline 的 535.8k 降至 22.4k（约一个数量级以上节省），但第 6 次考试的 DPAE 上升到 \(0.131\pm0.008\)。这表明“可信度估计”不仅是工程成本的主导项，也会实质性影响叙事性噪声的扩散与抑制，是值得作为推断子模块单独建模与校准的环节。
\paragraph{（iii）No-RAG：与 Baseline 差异有限，当前实现中 RAG 更像结构约束而非强信息增益。}
No-RAG（仅保留 self + class\_stats，不提供目标相关检索）在第 6 次考试的 DPAE 为 \(0.121\pm0.011\)，与 Baseline 的 \(0.124\pm0.009\) 区间重叠。更合理的解释是：当前 RAG 的主要贡献在于\textbf{硬约束信息可见范围}（防止越权），而证据设计的边际信息增益尚不足以在宏观误差上显著凸显；后续可通过更细粒度证据、或将来源置信度显式并入 \(\tau\) 的构造来提高可识别性。

\subsection{外部基线对比（External Baselines）：与纯预测范式的基准测试}
\label{subsec:external_baselines}

除上述针对模型机制的消融实验外，我们额外引入了一组\textbf{不依赖 LLM 交互}的外部基线。这一部分的目的是回答一个基础性问题：\emph{若将本任务简化为单纯的“期末（Epoch 6）班级排名预测”问题，而不考虑主观视野、叙事噪声与信任门控的模拟真实性，现有的启发式或监督学习方法表现如何？}

实验在 Social-Observed 样本集（419 人，12 个班级）上运行，最终考试（Epoch 6）的指标汇总于表~\ref{tab:external_baselines}。需注意，本文定义的 DPAE 为 \(\text{DPAE}=1-\rho\)（Spearman 相关系数的互补量）。因此，\textbf{随机排序的期望表现应接近 DPAE\(\approx 1\) 与 \(\rho\approx 0\)}，这为评估模型的有效性提供了下界参照。

\begin{table}[!htbp]
\centering
\small
\begin{tabular}{lcccc}
\toprule
方法 & DPAE$_6$ & Spearman $\rho_6$ & Acc@3$_6$ & Div$_6$ \\
\midrule
Random & $0.999\pm0.147$ & 0.001 & 0.028 & -- \\
Self-Only & $1.059\pm0.159$ & -0.059 & 0.111 & -- \\
Linear Regression & $0.838\pm0.173$ & 0.162 & 0.083 & -- \\
MLP Regressor & $0.839\pm0.121$ & 0.161 & 0.083 & -- \\
\midrule
SGC/GCN-like (1-hop) & $0.977\pm0.179$ & 0.023 & 0.111 & -- \\
SGC/GCN-like (2-hop) & $0.953\pm0.165$ & 0.047 & 0.111 & -- \\
GAT-like (1-hop) & $0.847\pm0.228$ & 0.153 & 0.139 & -- \\
DeGroot dynamics & $0.154\pm0.044$ & 0.846 & 0.361 & $5\times10^{-4}$ \\
\midrule
Ours (Rashomon Agents) & $\mathbf{0.124\pm0.009}$ & \textbf{0.876} & 0.278 & (see main exp.) \\
\bottomrule
\end{tabular}
\caption{外部基线对比实验：最终考试（Epoch 6）指标汇总。DPAE 定义为 \(1-\rho\)（越小越好）；Acc@3 为前三名识别率（越大越好）；Div$_6$ 表示 DeGroot 动力学的终态“意见多样性”（越小表示越接近共识塌缩）。}
\label{tab:external_baselines}
\end{table}

\paragraph{（i）Random / Self-Only：确立任务难度与交互必要性。}
Random 基线的 \(\rho\approx 0\) 及其极低的 Acc@3 证实了该任务无法通过随机猜测命中 Top-3。Self-Only（无社交融合，仅保留个体的固定先验）表现出轻微的负相关（\(\rho<0\)），这有力地说明：\textbf{在缺乏交互证据的情况下，仅凭个体的初始信念无法恢复真实的群体排序}。这一结果验证了实验设定的合理性，即社会信息的融合对于认知修正至关重要。

\paragraph{（ii）问卷特征回归（Linear/MLP）：静态特征的信息瓶颈。}
两类监督学习基线的表现均徘徊在 \(\rho\approx 0.16\) 左右，Acc@3 仅为 0.083。这一低性能的合理解释在于：问卷特征（多为独热编码的主观题）与“期末学业排名”之间不存在直接的强因果关联，且不同班级间存在显著的分布异质性。在样本量有限且特征稀疏的约束下，纯特征回归模型难以捕捉深层的动态关联，往往只能学到跨班级的混淆模式，无法达到基于动态交互证据所能实现的排序一致性。

\paragraph{（iii）图学习基线（SGC/GAT-like，1-hop）：静态数值聚合无法替代认知推理。}
为了确保比较的公平性并避免设置过弱的基线，我们构建的 Graph Learning 基线\textbf{严格遵循了与 Rashomon Agents 相同的罗生门视野约束}：即仅允许在由问卷好友块构建的一阶邻居子图上进行聚合。结果显示，SGC/GCN-like 的表现接近随机（\(\rho\approx 0.02\sim 0.05\)），GAT-like 也仅达到 \(\rho\approx 0.153\)，与线性回归相当，\textbf{未能从图结构中提取出有效的排序信号}。这一结果揭示了关键的方法论差异：在当前的“好友关系+问卷特征”设定下，传统的\textbf{静态数值聚合}机制失效了。一方面，好友边本身并不直接承载学业强弱的数值信号；另一方面，简单的 1-hop 聚合在稀疏、噪声大且异质的图结构中极易导致信息过平滑（Over-smoothing）。相比之下，Rashomon Agents 通过\textbf{叙事性交互}与\textbf{信任门控}机制，显式地建模了“关系—证据—更新”的语义过程，从而在同等的局部视野限制下，实现了显著更优的排序恢复能力。

\paragraph{（iv）DeGroot（意见动力学）：共识塌缩与同质化困境。}
DeGroot 动力学虽然在预测精度上表现尚可（DPAE$_6 \approx 0.154$, \(\rho_6\approx 0.846\)），显示了线性共识扩散在纯预测任务上的竞争力，但其终态多样性指标 Div$_6\approx 5\times 10^{-4}$ 揭示了其致命缺陷：群体意见几乎完全塌缩为单一共识。更关键的是，DeGroot 的指标在 Epoch 1 与 Epoch 6 间几乎无变化，呈现出“\textbf{快速收敛至固定点}”的特征。这意味着它无法复现我们所关注的\textbf{误判随时间累积}与\textbf{自信地错误}等复杂社会现象，也无法维持接近真实人类社会的“\textbf{多样化错误}”。因此，DeGroot 仅可视为一个理想化的无偏共识参照系。与之相对，我们的 Agent 模型在保持高排序一致性的同时，成功保留了信念结构的异质性，从而能解释现实中“集体误判与个体分歧共存”的现象。

为进一步剖析 DeGroot 的动力学特性，我们在 Epoch 6 扫描了不同的迭代步数（steps），结果详见表~\ref{tab:degroot_sweep_epoch6} 与图~\ref{fig:degroot_sweep_epoch6}。我们观察到两个稳健的规律：
\begin{itemize}
    \item \textbf{过度收敛导致预测退化}：随着步数增加（steps=1 $\to$ 100），\(\rho\) 值反而从 0.952 降至 0.825。这直观地表明，线性的平均化操作逐渐“洗去”了个体私有信号中的有效差异，导致群体估计向低信息的“均匀共识”退化。
    \item \textbf{持续极低的多样性}：无论迭代步数如何，系统均快速塌缩至近共识状态（Div $\sim 10^{-4}$）。这与本文旨在探究的、由叙事驱动和认知偏差主导的路径依赖现象形成了鲜明对比。
\end{itemize}

\begin{table}[!htbp]
\centering
\small
\begin{tabular}{rcccc}
\toprule
steps & DPAE$_6$ & Spearman $\rho_6$ & Acc@3$_6$ & Div$_6$ \\
\midrule
1 & $0.048\pm0.015$ & 0.952 & 0.528 & 0.000727 \\
5 & $0.103\pm0.035$ & 0.897 & 0.444 & 0.000472 \\
30 & $0.154\pm0.044$ & 0.846 & 0.361 & 0.000513 \\
100 & $0.175\pm0.058$ & 0.825 & 0.333 & 0.000542 \\
\bottomrule
\end{tabular}
\caption{DeGroot 动力学步数扫描（Epoch 6）。Div$_6$ 为终态意见多样性（平均方差），数值越小表示越接近共识塌缩。完整数据见 \texttt{output/degroot\_sweep/degroot\_sweep\_epoch6.csv}。}
\label{tab:degroot_sweep_epoch6}
\end{table}

\begin{figure}[!htbp]
\centering
\includegraphics[width=0.95\linewidth]{output/degroot_sweep/degroot_sweep_epoch6.png}
\caption{DeGroot 步数扫描可视化（Epoch 6）：左图展示 DPAE 与 Spearman 相关系数随迭代步数（对数坐标）的变化趋势；右图展示终态多样性 Div$_6$。结果表明，随着线性共识迭代的加深，有效的排序信息被逐渐“平滑”掉，且系统快速陷入共识塌缩状态。}
\label{fig:degroot_sweep_epoch6}
\end{figure}
\FloatBarrier

\subsection{单种子端到端运行（Pilot sanity check）：现象可视化与班级异质性}
为验证工程端到端可用性并提供直观轨迹，我们在真实数据上完成一次完整模拟（seed=42），开启 100 并行线程以加速 LLM 调用；完整 artefacts 保存在 \texttt{output/01-22-00-31/}。该单次运行的数值与趋势与多随机种子主结果保持一致，并额外揭示了两个可用于论文叙事的现象：\textbf{（a）误判随时间累积的可视化轨迹}，以及\textbf{（b）“自信地错误”与班级异质性}。

\paragraph{误判随时间累积（单种子示例）。}
表~\ref{tab:results}（由 \texttt{output/01-22-00-31/figures/results\_table.tex} 插入）与图~\ref{fig:dpae_trend} 显示：Spearman 相关 \(\rho\) 随 epoch 单调下降（0.9279 \(\rightarrow\) 0.8653），对应 DPAE 单调上升（0.0721 \(\rightarrow\) 0.1347），呈现宏观误判的累积性偏离。

\input{output/01-22-00-31/figures/results_table.tex}
\FloatBarrier

\paragraph{“自信地错误”（单种子示例）。}
图~\ref{fig:uncertainty_trend} 将平均信念不确定性与 Top-$k$ 识别率并置：平均 \(\sigma\) 下降（0.2715 \(\rightarrow\) 0.2359），但 Top-3 学霸识别率同时从 0.69 下降到 0.22（Top-5 亦下降），体现出 ``confidently wrong'' 的背离现象，即群体主观上更确定、客观上更不准。

\paragraph{班级间异质性（单种子示例）。}
图~\ref{fig:class_comparison} 与图~\ref{fig:dpae_by_class} 表明不同班级的最终 DPAE 在约 0.07--0.21 区间内波动，且随时间上升的斜率与幅度存在差异。这提示网络结构与人群心理构成会调制宏观误判强度，也为后续“噪声放大器”检验提供了统计可识别的方差来源。

\begin{figure}[!htbp]
\centering
\includegraphics[width=0.95\linewidth]{output/01-22-00-31/figures/dpae_trend.png}
\caption{聚合层面的 DPAE（左轴）与 Spearman 相关（右轴）随时间变化。DPAE 单调上升而相关性下降，呈现宏观误判的累积性偏离。}
\label{fig:dpae_trend}
\end{figure}

\begin{figure}[!htbp]
\centering
\includegraphics[width=0.95\linewidth]{output/01-22-00-31/figures/uncertainty_trend.png}
\caption{平均信念不确定性与 Top-$k$ 学霸识别率随时间变化。不确定性下降但 Top-$k$ 同时下降，提示“收敛到错误共识”的动力学现象。}
\label{fig:uncertainty_trend}
\end{figure}

\begin{figure}[!htbp]
\centering
\includegraphics[width=0.95\linewidth]{output/01-22-00-31/figures/class_comparison.png}
\caption{各班级最终 DPAE（Epoch 6）对比。不同班级的误判强度差异显著，提示网络结构与人群心理构成对宏观误判具有调制作用。}
\label{fig:class_comparison}
\end{figure}

\begin{figure}[!htbp]
\centering
\includegraphics[width=0.95\linewidth]{output/01-22-00-31/figures/dpae_by_class.png}
\caption{各班级 DPAE 随时间的轨迹。多数班级呈现随时间上升趋势，但斜率与幅度存在差异。}
\label{fig:dpae_by_class}
\end{figure}

\subsection{可视化补充：跨 epoch 轨迹、终点对比与配对差异}
图~\ref{fig:uai_traj} 以跨种子均值曲线与标准差带展示 DPAE 与 Spearman 的跨 epoch 轨迹；图~\ref{fig:uai_traj_ci} 进一步用 bootstrap 95\% CI 更稳健地表达跨种子不确定性。图~\ref{fig:uai_finalbars} 将最终 epoch（第 6 次考试）的 DPAE、Spearman 与 Top-$k$ 指标并置对比，便于快速判断消融差异的幅度与稳定性。最后，图~\ref{fig:uai_paired} 将消融相对 Baseline 的净效应以配对差异形式呈现，直接回答“去掉某机制后的性能变化及其不确定性”。
\FloatBarrier

\section{统计检验与复现协议}
\subsection{统计量与不确定性报告}
所有主结论将以\textbf{跨随机种子}（seed）与\textbf{跨班级}（class）的双层统计方式报告。具体地，我们先对每个 seed 在每个班级上计算指标（DPAE、Unc、Strat、RecoveryTime/Area 等），再在班级内取均值形成 \(\hat{m}_s\)，最后在 seeds 上报告均值与置信区间。

我们将采用以下两种互补的区间估计：
\begin{enumerate}
    \item \textbf{Bootstrap 置信区间}：对 seeds 维度重采样（例如 20,000 次）得到 95\% CI；
    \item \textbf{分层置换检验}：以班级为分层单位置换标签，检验“噪声放大器（emotionally\_volatile）”效应或 no-RAG/no-Subjective-Graph 的差异是否显著。
\end{enumerate}
这使得我们在样本规模不大（12 班）时仍能给出稳健统计推断。

\subsection{随机性与缓存策略}
为保证复现，我们固定以下随机源：
\[
\texttt{seed}\ \Rightarrow\ \text{交互对象采样、批处理目标采样、噪声注入的伪随机序列。}
\]
同时，我们对 LLM 调用采用“输入哈希 \(\to\) 输出缓存”的确定性缓存，以减少波动与成本；所有缓存与日志将随实验 artefacts 一并保存。

\section{伦理与隐私}
该数据集包含学生心理与社交信息，属于高度敏感数据。实验报告与论文中仅以匿名 id 呈现；RAG 检索内容严格限制在“同班聚合统计或匿名片段”，避免直接输出任何可识别个人信息。

\section{预期贡献}
\begin{itemize}
    \item \textbf{方法贡献}：提出 Rashomon set agents 的形式化（主观图集合）与 RAG 受限视野下的动态信念传播框架。
    \item \textbf{实证贡献}：在稀有的“心理问卷 + 交互关系 + 多次考试”真实数据上，验证微观主观交互是否足以涌现宏观学业分化与集体误判。
    \item \textbf{机制贡献}：量化情绪波动/焦虑等微观变量在不确定性传播中的角色，给出“噪声放大器”与级联失败的可检验指标体系。
\end{itemize}

\appendix

\section{LLM Agent 提示词与结构化输出协议（Prompts \& Schemas）}
\subsection{能力评估（批处理）提示词}
系统提示词（System）：
\begin{verbatim}
你是班级社交环境中的学生智能体。你只能依据给定的【检索证据】与【自身性格/焦虑】做出判断。

你的任务是评估若干同学的学业能力。请输出严格 JSON，不要输出任何额外文本。

输出格式:
{
  "evaluations": [
    {
      "target_id": <目标学生ID>,
      "ability_score": <0-1之间的能力评分，1表示最强>,
      "uncertainty": <0-1之间的不确定性，1表示完全不确定>,
      "rationale": "<≤40字的评价理由>"
    }
  ]
}

注意:
- ability_score 应基于你对该同学学业表现的主观感知
- uncertainty 反映你对这个评估的确信程度
- 如果你对某同学了解很少，uncertainty 应该较高
- 理由要简洁，不超过40个字
\end{verbatim}

用户提示词（User，含可检索证据与信念摘要）：
\begin{verbatim}
【你的信息】
- 学生ID: {agent_id}
- 性格: {personality}
- 焦虑水平: {worry_level}

【检索证据】
{rag_context}

【当前信念摘要】
{belief_summary}

【任务】
请评估以下同学的学业能力: {target_ids}

请只输出JSON，不要有其他内容。
\end{verbatim}

\subsection{可信度/信任门控提示词（LLM-Trust）}
系统提示词（System）：
\begin{verbatim}
你是信息可信度评估器。根据发送者/接收者性格、上下文一致性，评估收到的评价信息的可信度。

输出格式（只输出JSON）:
{
  "trust_weight": <0-1之间的信任度，1表示完全信任>,
  "uncertainty": <0-1之间的不确定性>
}

评估依据:
- 发送者的性格特点（情绪稳定的人可能更可靠）
- 评价与你当前认知的一致性
- 评价的不确定性（高不确定性的评价可能需要打折）
\end{verbatim}

用户提示词（User）：
\begin{verbatim}
【你的信息】
- 学生ID: {receiver_id}
- 性格: {receiver_personality}

【收到的评价】
- 发送者ID: {sender_id}
- 发送者性格: {sender_personality}
- 关于: 学生{target_id}
- 能力评分: {ability_score}
- 不确定性: {sender_uncertainty}
- 理由: {rationale}

【你对学生{target_id}的当前认知】
- 能力均值: {current_mu:.2f}
- 不确定性: {current_sigma:.2f}

请只输出JSON，不要有其他内容。
\end{verbatim}

\end{document}